\chapter{Reprodução}\label{apen:apendiceA}


Neste capítulo, é disponibilizado todas as informações necessárias para reproduzir os resultados.

\section{Hiperparâmetros do OPTICS}

    \begin{itemize}
        \item Aniso e Blobs: min\_samples=$7$, xi=$0.1$, min\_cluster\_size=$0.2$;
        \item Varied: min\_samples=$7$, xi=$0.01$, min\_cluster\_size=0.2;
        \item Moons: min\_samples=$7$, xi=$0.1$, min\_cluster\_size=0.1;
        \item Circles: min\_samples: $7$, xi=$0.08$, min\_cluster\_size=0.1;
        \item No Structure: min\_samples=$7$, xi=$0.05$, min\_cluster\_size=$0.1$.
    \end{itemize}  

\section{Hiperparâmetros para o $\beta^{4}$-\gls{IRT}}

Nesta seção, uma visão geral dos parâmetros para configuração do modelo $\beta^{4}$-\gls{IRT} é disponibilizada. Abaixo, temos os parâmetros utilizados:
\begin{itemize}
    \item \textbf{Learning Rate:} $0.1$
    \item \textbf{Epochs:} $10000$
    \item \textbf{Number of Initializations:} $1000$
    \item \textbf{Random Seed:} $1$
    \item \textbf{Tolerance:} $1.0 \times 10^{-8}$
    \item \textbf{Set Priors:} \texttt{True}
\end{itemize}

Estas configurações influenciam no processo de otimização, criterio de convergência e estrategia de inicialização para o $\beta^4$-\gls{IRT}. É imporante ajustar estes parâmetros baseado nestas características  especificas para certificação de reprodutibilidade.


\section{Configurações dos conjuntos de dados sintéticos}

Todos os conjuntos de dados sintéticos foram criados com tamanho de amostra de $500$, exceto o conjunto No Structure, que foi gerado com tamanho de amostra de $1000$. O estado aleatório foi consistentemente definido como $0$ para todos os conjuntos de dados utilizando a biblioteca \textsf{NumPy} \cite{numpy}, com exceção dos conjuntos No Structure e Varied, nos quais foi utilizado o estado aleatório $170$, e do conjunto Blobs, no qual o estado aleatório foi especificamente definido como $8$.

\begin{itemize}
    \item \textbf{Aniso:}
    \begin{itemize}
        \item Método: O conjunto sintético Aniso foi gerado utilizando o método \textsf{make\_blobs} do Scikit-Learn, permitindo a criação de agrupamentos anisotrópicos.
        \item Matrizes de rotação: [$0.6$, $-0.6$], [$-0.4$, $0.8$]
    \end{itemize}
    \item \textbf{Noisy Circles:}
    \begin{itemize}
        \item Método: O conjunto Noisy Circles foi gerado utilizando o método \textsf{make\_circles} do Scikit-Learn, introduzindo ruído em agrupamentos circulares.
        \item Fator de ruido: $0.5$
        \item Ruído: $0.05$
    \end{itemize}
    \item \textbf{Noisy Moons:}
    \begin{itemize}
        \item Método: O conjunto Noisy Moons foi gerado utilizando o método \textsf{make\_circles} do Scikit-Learn, resultando em agrupamentos em formato de lua com adição de ruído.
        \item Ruído: $0.05$
    \end{itemize}
    \item \textbf{Blobs:}
    \begin{itemize}
        \item Método: O conjunto Blobs foi gerado utilizando o método \textsf{make\_blobs} do Scikit-Learn, produzindo agrupamentos Gaussianos com centros e desvios padrão especificados.
        \item \textit{Random state}: $8$
    \end{itemize}
    \item \textbf{Varied:}
    \begin{itemize}
        \item Método: O conjunto Varied foi gerado utilizando o método \textsf{make\_blobs} do Scikit-Learn, criando agrupamentos Gaussianos escalados com diferentes desvios padrão entre os clusters.
        \item Desvios padrão dos clusters: $1.0$, $2.5$ e $0.5$
    \end{itemize}
    \item \textbf{No Structure:}
    \begin{itemize}
        \item Método: O conjunto No Structure foi gerado com tamanho de amostra e número de atributos específicos, com o objetivo de representar um conjunto de pontos distribuídos aleatoriamente, sem qualquer estrutura inerente.
        \item Tamanho da amostra: $1000$
        \item Número de atributos: $2$
    \end{itemize}
\end{itemize}

Esses parâmetros de geração permitem a criação reprodutível dos conjuntos de dados sintéticos, facilitando a experimentação consistente e a avaliação de algoritmos de agrupamento.

\section{Geração de partições aleatórias}

O Algoritmo \ref{alg:random:partitions} descreve as etapas para gerar partições aleatórias e incorporá-las ao conjunto de modelos.

\begin{algorithm}
\caption{Incorporação de Partições Aleatórias}
\begin{algorithmic}
% \State Gerar uma matriz aleatória, com $N$ instâncias e $P$ partições.
\State Para cada partição gerada, fixar um \textit{Random State} $p+1$, onde $p = 0, 1, 2, \dots, P$.

\For{cada conjunto de dados}
    \For{$p = 0, 1, \dots, P$}
    \State Ajustar cada modelo de agrupamento e agrupar as instâncias.
    
    %\State Adição de $p$ partições para cada instância a um dos $K$ clusters.

    \State Atribuir aleatoriamente cada instância a um dos $K$ grupos, onde $K$ é o número conhecido de clusters no conjunto de dados, ou $K=10$ para o conjunto No Structure.
    
    \State Calcular a matriz de respostas.
    
    \State Calcular os modelos artificiais (\textit{Best} e \textit{Average response} (as partições aleatórias não são consideradas para determinar suas respostas).
    
    \State Ajustar o modelo $\beta^{4}$-\gls{IRT} com a matriz de respostas ($p_{ij}$).
    \EndFor
\EndFor
\end{algorithmic}
\label{alg:random:partitions}
\end{algorithm}